\documentclass[11pt, a4paper]{article}

\usepackage[margin=2.5cm]{geometry}
\usepackage[T1]{fontenc}
\usepackage[utf8]{inputenc}
\usepackage{titlesec}
\usepackage{enumitem}
\usepackage{booktabs}
\usepackage{tabularx}
\usepackage{longtable}
\usepackage{array}
\usepackage{parskip}
\usepackage{fancyhdr}
\usepackage{xcolor}
\usepackage{mdframed}
\usepackage{listings}
\usepackage{hyperref}
\usepackage{amsmath}
\usepackage{tocloft}
\usepackage{multicol}

% Colors
\definecolor{rulecolor}{RGB}{40,40,40}
\definecolor{phasecolor}{RGB}{20,80,140}
\definecolor{warncolor}{RGB}{160,40,40}
\definecolor{codebackground}{RGB}{246,246,246}
\definecolor{codekeyword}{RGB}{0,80,160}
\definecolor{codestring}{RGB}{140,40,0}
\definecolor{codecomment}{RGB}{80,110,80}
\definecolor{tiercolor}{RGB}{60,120,60}
\definecolor{fillcolor}{RGB}{100,60,140}

\titleformat{\section}{\large\bfseries\color{phasecolor}}{}{0em}{}[\titlerule]
\titleformat{\subsection}{\normalsize\bfseries\color{phasecolor!70!black}}{}{0em}{}
\titleformat{\subsubsection}{\normalsize\itshape}{}{0em}{}

\lstset{
  backgroundcolor=\color{codebackground},
  basicstyle=\ttfamily\footnotesize,
  keywordstyle=\color{codekeyword}\bfseries,
  stringstyle=\color{codestring},
  commentstyle=\color{codecomment}\itshape,
  breaklines=true,
  frame=single,
  framerule=0.4pt,
  rulecolor=\color{rulecolor!30},
  xleftmargin=8pt,
  xrightmargin=8pt,
  aboveskip=8pt,
  belowskip=8pt,
  showstringspaces=false,
  language=C++
}

\pagestyle{fancy}
\fancyhf{}
\renewcommand{\headrulewidth}{0.4pt}
\fancyhead[L]{\small\texttt{UFX\_AUTO\_2} --- Database Fillout Phases 6--10}
\fancyhead[R]{\small VSEPR-SIM / v5 beta9}
\fancyfoot[C]{\small\thepage}

\newmdenv[
  linecolor=warncolor,
  linewidth=1pt,
  leftmargin=0pt, rightmargin=0pt,
  innerleftmargin=10pt, innerrightmargin=10pt,
  innertopmargin=8pt, innerbottommargin=8pt,
  backgroundcolor=warncolor!4
]{rulebox}

\newmdenv[
  linecolor=phasecolor,
  linewidth=1pt,
  leftmargin=0pt, rightmargin=0pt,
  innerleftmargin=10pt, innerrightmargin=10pt,
  innertopmargin=8pt, innerbottommargin=8pt,
  backgroundcolor=phasecolor!4
]{phasebox}

\newmdenv[
  linecolor=tiercolor,
  linewidth=1pt,
  leftmargin=0pt, rightmargin=0pt,
  innerleftmargin=10pt, innerrightmargin=10pt,
  innertopmargin=8pt, innerbottommargin=8pt,
  backgroundcolor=tiercolor!4
]{tierbox}

\newmdenv[
  linecolor=fillcolor,
  linewidth=1pt,
  leftmargin=0pt, rightmargin=0pt,
  innerleftmargin=10pt, innerrightmargin=10pt,
  innertopmargin=8pt, innerbottommargin=8pt,
  backgroundcolor=fillcolor!4
]{fillbox}

\renewcommand{\cftsecleader}{\cftdotfill{\cftdotsep}}
\setlength{\cftbeforesecskip}{2pt}

% Relax overfull box tolerance for long monospace lines in listings/tables.
\sloppy
\setlength{\emergencystretch}{3em}

\begin{document}

% =========================================================
% TITLE
% =========================================================
\begin{center}
  {\LARGE\bfseries UFX\_AUTO\_2}\\[4pt]
  {\large\bfseries Database Fillout Work Order: Phases 6--10}\\[6pt]
  {\normalsize VSEPR-SIM $\cdot$ v5 beta9 target}\\[3pt]
  {\small\texttt{Prerequisite: tiny/medium continual run stable (Phases 2--5 complete)}}\\[3pt]
  {\small\textit{Phase 2 fills the skeleton. Phases 6--10 give it something to say.}}
\end{center}

\vspace{4pt}
\hrule height 0.8pt
\vspace{10pt}

\tableofcontents
\newpage

% =========================================================
\section{Purpose and Context}
% =========================================================

Phase 2 generated the atomistic identity and force-field spine.  It proved the
generation loop, the validation path, the web cache, and the continual harness.
After Phase 2, the database holds approximately 100 unique \texttt{material\_key}
records before the axis space saturates, because only five axes vary:
\textit{element}, \textit{oxidation\_state}, \textit{coordination},
\textit{geometry}, and \textit{phase}.

\begin{rulebox}
Phases 6--10 are not a rewrite.  They are additive.  Each phase extends the
existing \texttt{UFXMaterialRecord}, expands the generation axis space, adds one
or more generator classes, and wires one or more new CLI sub-commands into the
existing \texttt{ufx auto2} dispatch chain.  The harness loop does not change.
\end{rulebox}

Every phase follows the same template as Phase 2:

\begin{enumerate}[leftmargin=*, itemsep=2pt]
  \item Extend \texttt{AxisConfig} with new sampling axes
  \item Implement a \texttt{*Generator} class that fills the corresponding record block
  \item Extend \texttt{LocalSanityChecker} with tier-appropriate checks
  \item Insert block data into \texttt{property\_values} with correct \texttt{block\_tier}/\texttt{block\_name}
  \item Add a CLI sub-command: \texttt{ufx auto2 fill-<tier>}
  \item Wire into the continual harness cycle after \texttt{randomfill}
\end{enumerate}

\subsection{Axis space growth by phase}

\begin{center}
\begin{tabular}{clrr}
\toprule
\textbf{Phase} & \textbf{Block} & \textbf{New axes} & \textbf{Approx.\ unique key growth} \\
\midrule
2  & Identity + Force-field          & 5   & $\sim$100 \\
6  & Molecular descriptors           & +4  & $\times$10 \\
7  & Thermo + EOS                    & +6  & $\times$20 \\
8  & Crystal / solid-state           & +5  & $\times$30 \\
9  & Transport + Mechanics           & +6  & $\times$25 \\
10 & Meta-score active scheduler     & n/a & coverage-driven steering \\
\bottomrule
\end{tabular}
\end{center}

\begin{phasebox}
Phase 10 does not add new record axes.  It steers the sampler toward under-covered
regions of the combined axis space defined by Phases 2--9.  Without it, the
continual loop will over-generate common elements at standard conditions and
under-represent actinides, high-pressure phases, and exotic geometries.
\end{phasebox}

% =========================================================
\section{Shared Infrastructure (Review Before Phase 6)}
% =========================================================

These components are already built and must be understood before extending them.

\subsection{AxisConfig}

\texttt{include/ufx\_auto2/axis\_config.hpp} defines the sampling envelope.
\texttt{AxisConfig::default\_phase2()} populates five axis vectors.  Phases 6--9
each call a new static builder, \texttt{AxisConfig::extend\_phase\textit{N}()}, that
adds axes to an existing config rather than replacing it.

\subsection{AxisSample}

\texttt{include/ufx\_auto2/axis\_config.hpp} also defines \texttt{AxisSample}.
New axes are added as fields here.  Use \texttt{std::optional<>} for axes that
are not populated in all phases.

\begin{lstlisting}
struct AxisSample {
	// Phase 2 (mandatory)
	std::string element_family;
	std::string element;
	int         oxidation_state;
	int         coordination;
	std::string geometry;
	std::string phase;
	double      temperature_K;
	double      pressure_atm;
	std::string target_property;
	int         sample_index;
	uint64_t    seed_used;

	// Phase 6 -- molecular
	std::optional<std::string> canonical_smiles;
	std::optional<int>         heavy_atom_count;
	std::optional<double>      molecular_weight;

	// Phase 7 -- thermo
	std::optional<double> delta_Hf_kJ_mol;
	std::optional<double> cp_298_J_mol_K;
	std::optional<std::string> eos_model;

	// Phase 8 -- crystal
	std::optional<std::string> space_group;
	std::optional<double>      lattice_a_angstrom;
	std::optional<std::string> structure_source;

	// Phase 9 -- transport/mechanics
	std::optional<double>      strain_rate;
	std::optional<std::string> test_environment;
};
\end{lstlisting}

\subsection{property\_values insert pattern}

Every new block value is inserted via \texttt{ufx\_insert\_property\_value()} using
the correct \texttt{block\_tier} and \texttt{block\_name} from the table below.
State columns (\texttt{temperature\_K}, \texttt{pressure\_Pa}, \texttt{phase})
are always populated.  Do not leave them at default without intent.

\begin{center}
\begin{tabular}{clll}
\toprule
\textbf{Tier} & \textbf{block\_name} & \textbf{Phase} & \textbf{Example property\_name} \\
\midrule
0 & \texttt{identity}    & 2  & \texttt{coordination\_number}, \texttt{geometry\_tag} \\
1 & \texttt{force\_field} & 2  & \texttt{r1}, \texttt{D1}, \texttt{theta0} \\
2 & \texttt{molecular}    & 6  & \texttt{molecular\_weight}, \texttt{inchikey} \\
3 & \texttt{thermo}       & 7  & \texttt{delta\_Hf\_kJ\_mol}, \texttt{cp\_298} \\
3 & \texttt{eos}          & 7  & \texttt{density\_kg\_m3}, \texttt{viscosity\_Pa\_s} \\
4 & \texttt{crystal}      & 8  & \texttt{space\_group}, \texttt{density\_g\_cm3} \\
5 & \texttt{mechanical}   & 9  & \texttt{E\_eff\_GPa}, \texttt{yield\_MPa} \\
5 & \texttt{transport}    & 9  & \texttt{k\_eff\_W\_mK}, \texttt{diffusivity\_m2\_s} \\
9 & \texttt{meta}         & 10 & \texttt{weirdness}, \texttt{coverage}, \texttt{composite\_score} \\
\bottomrule
\end{tabular}
\end{center}

% =========================================================
\section{Phase 6 --- Molecular Descriptor Importer (Tier 2)}
% =========================================================

\begin{tierbox}
\textbf{Goal:} attach \texttt{MolecularDescriptorBlock} data to existing generated records,
validate via PubChem PUG REST, and expand the axis space with molecular-formula
and structural dimensions.
\end{tierbox}

\subsection{Why Phase 6 first}

Molecular descriptors are cheap to generate (formula arithmetic) and cheap to
validate (one PubChem CID lookup).  They act as the first meaningful
external-reference check after the Phase 5 web cache is warm.  They also unlock
the Tier 2 sanity filter: a generated \texttt{CH3Br} entry with the wrong
molecular weight should fail before thermo data is ever computed.

\subsection{New axis dimensions}

\begin{itemize}[leftmargin=*, itemsep=2pt]
  \item \textbf{canonical\_smiles} --- sampled from a SMILES vocabulary per element family
  \item \textbf{heavy\_atom\_count} --- integer 1--20
  \item \textbf{hbond\_donor\_count} --- integer 0--6
  \item \textbf{hbond\_acceptor\_count} --- integer 0--8
\end{itemize}

SMILES are not random strings.  They are drawn from a curated per-element table
that encodes known coordination/oxidation combinations.  This table lives in
\texttt{include/ufx\_auto2/smiles\_vocab.hpp}.

\subsection{Generator class: MolecularDescriptorGenerator}

File: \texttt{src/ufx\_auto2/molecular\_descriptor\_generator.cpp}

\begin{lstlisting}
class MolecularDescriptorGenerator {
public:
	MolecularDescriptorGenerator();

	// Fill rec.molecular from axis sample + periodic table data.
	// Does NOT call the network.  All values are deterministic.
	void fill(UFXMaterialRecord& rec, const AxisSample& s) const;

	// Compute molecular weight from elements + stoichiometry.
	static double molecular_weight_from_formula(const std::string& formula);

	// Build a canonical InChIKey approximation from element + oxidation.
	static std::string inchikey_stub(const std::string& element,
									 int ox_state,
									 int coordination);
private:
	// Periodic table atomic weights (local; no external lookup).
	static double atomic_weight_(const std::string& sym) noexcept;
};
\end{lstlisting}

The generator fills \texttt{molecular\_formula}, \texttt{molecular\_weight},
\texttt{exact\_mass}, \texttt{heavy\_atom\_count}, and a stub \texttt{inchikey}.
It does not attempt SMILES generation from scratch --- that is a Phase 6b concern.
For Phase 6, SMILES is either sampled from the vocabulary table or left empty.

\subsection{LocalSanityChecker extension}

Add the following Tier 2 checks to \texttt{LocalSanityChecker::check()}:

\begin{itemize}[leftmargin=*, itemsep=2pt]
  \item \texttt{molecular.molecular\_weight > 0} --- must be positive
  \item \texttt{molecular.heavy\_atom\_count >= 1} --- at least one heavy atom
  \item molecular weight is consistent with formula (within 5\% tolerance)
  \item \texttt{inchikey} is non-empty if identity block is fully populated
\end{itemize}

Failure mode tag: \texttt{molecular\_descriptor\_invalid}.

\subsection{property\_values rows written by fill-molecular}

\begin{center}
\begin{tabular}{llll}
\toprule
\textbf{property\_name} & \textbf{block\_tier} & \textbf{units} & \textbf{source} \\
\midrule
\texttt{molecular\_weight}       & 2 & g/mol    & \texttt{periodic\_table\_arithmetic} \\
\texttt{exact\_mass}             & 2 & g/mol    & \texttt{periodic\_table\_arithmetic} \\
\texttt{heavy\_atom\_count}      & 2 & count    & \texttt{formula\_parse} \\
\texttt{hbond\_donor\_count}     & 2 & count    & \texttt{smiles\_count} \\
\texttt{hbond\_acceptor\_count}  & 2 & count    & \texttt{smiles\_count} \\
\texttt{inchikey}                & 2 & string   & \texttt{inchi\_stub} \\
\bottomrule
\end{tabular}
\end{center}

\subsection{Web cross-check (validate-web extension)}

The Phase 5 PubChem fetcher already retrieves \texttt{molecular\_weight} and
\texttt{inchikey}.  Phase 6 adds a comparison step inside
\texttt{WebValidator::validate\_batch()}:

\begin{enumerate}[leftmargin=*, itemsep=2pt]
  \item Lookup PubChem CID by formula or element symbol
  \item Compare \texttt{MolecularWeight} response field vs.\ stored
		\texttt{property\_values} row (block=\texttt{molecular}, property=\texttt{molecular\_weight})
  \item If relative error $> 2\%$: write \texttt{validation\_records} row with
		\texttt{passed=0}, note \texttt{molecular\_weight\_mismatch}
  \item If CID not found: write \texttt{passed=1}, note \texttt{web\_missing}, confidence capped 0.85
\end{enumerate}

\subsection{CLI sub-command}

\begin{lstlisting}[language={}]
vsepr-ufx ufx auto2 fill-molecular --batch 500 --db ufx_auto2.sqlite
vsepr-ufx ufx auto2 fill-molecular --batch 500 --db ufx_auto2.sqlite --verbose
\end{lstlisting}

Processes records where \texttt{block\_name='molecular'} rows are absent in
\texttt{property\_values}.  Writes molecular descriptor rows, then queues
for web cross-check on the next \texttt{validate-web} pass.

\subsection{Continual harness integration}

Insert after \texttt{randomfill}, before \texttt{validate}:

\begin{lstlisting}[language={}]
# Cycle order with Phase 6:
randomfill  --count $BatchSize
fill-molecular --batch $BatchSize
validate    --batch $ValidateBatch
validate-web --batch $WebBatch
promote     --min-score $MinScore
audit
\end{lstlisting}

\subsection{Audit success criteria after Phase 6}

\begin{itemize}[leftmargin=*, itemsep=2pt]
  \item \texttt{Total property values} nonzero and growing each cycle
  \item \texttt{Missing provenance = 0} (fill-molecular writes provenance rows)
  \item At least one \texttt{molecular\_weight\_mismatch} in rejected reasons
		(confirms the cross-check fires)
  \item Web validation: \texttt{web missing} acceptable, \texttt{web pass} nonzero
\end{itemize}

% =========================================================
\section{Phase 7 --- Thermo + EOS Blocks (Tier 3)}
% =========================================================

\begin{tierbox}
\textbf{Goal:} populate \texttt{ThermoBlock} and \texttt{EOSBlock} for generated records
using NIST WebBook as the primary external reference, periodic trend heuristics
as the generation fallback, and T/P-aware property curves stored in
\texttt{property\_values}.
\end{tierbox}

\subsection{Axis expansion}

\begin{itemize}[leftmargin=*, itemsep=2pt]
  \item \textbf{delta\_Hf\_kJ\_mol} --- sampled from heuristic periodic trend ranges
  \item \textbf{cp\_298\_J\_mol\_K} --- seed value; full curve generated from Shomate heuristic
  \item \textbf{eos\_model} --- one of: \texttt{ideal}, \texttt{vdW}, \texttt{PR}, \texttt{tabular}
  \item \textbf{temperature\_range} --- (T\_min, T\_max) bracket for the property curve
  \item \textbf{phase\_transition\_flag} --- boolean; indicates whether the record spans a
		phase boundary in the sample range
  \item \textbf{vapor\_pressure\_seed} --- log10 seed value at 298 K
\end{itemize}

\subsection{Generator class: ThermoGenerator}

File: \texttt{src/ufx\_auto2/thermo\_generator.cpp}

\begin{lstlisting}
class ThermoGenerator {
public:
	ThermoGenerator();

	// Populate rec.thermo using:
	//   1. Element-family heuristic ranges for Hf, Cp, Tb, Tm
	//   2. Shomate polynomial estimate (6 coefficients from group/period)
	//   3. EOS parameter estimates from critical property correlations
	void fill(UFXMaterialRecord& rec, const AxisSample& s) const;

	// Build a Cp(T) curve from Shomate heuristic in J/(mol K).
	static PropertyCurve estimate_cp_shomate(int Z,
											 double T_min_K,
											 double T_max_K,
											 int n_points = 20);

	// Estimate critical properties via Lee-Kesler correlations.
	static void estimate_critical_properties(UFXMaterialRecord& rec) noexcept;

private:
	// Group/period offsets for Cp seed values.
	static double cp_seed_(int period, const std::string& family) noexcept;
	static double Hf_seed_(int Z, int ox_state)                   noexcept;
};
\end{lstlisting}

\subsection{Shomate heuristic outline}

The Shomate polynomial is $C_p(T) = A + BT + CT^2 + DT^3 + E/T^2$.
For Phase 7, coefficients A--E are seeded from a lookup by (period, element family):

\begin{itemize}[leftmargin=*, itemsep=2pt]
  \item A$_0$ from monatomic Dulong-Petit limit: $3R \approx 24.9$ J/(mol\,K)
  \item B, C from polynomial regression offsets per period (stored in a 7-row table)
  \item D, E are zero for Phase 7 (first-order heuristic only)
\end{itemize}

This is explicitly labelled \texttt{method\_tag = "shomate\_heuristic\_phase7"} and
\texttt{confidence = 0.35} in the provenance row.  No false authority.

\subsection{EOS model selection rules}

\begin{center}
\begin{tabular}{ll}
\toprule
\textbf{Condition} & \textbf{EOS assigned} \\
\midrule
phase = gas, $T > 2T_c$       & ideal \\
phase = gas, $T \leq 2T_c$    & Peng-Robinson \\
phase = liquid                 & Peng-Robinson \\
phase = solid                  & tabular (density only) \\
phase = molten\_salt            & Peng-Robinson + salt correction flag \\
\bottomrule
\end{tabular}
\end{center}

\subsection{NIST WebBook fetcher extension}

The Phase 5 \texttt{NISTFetcher} is extended with two new query modes:

\begin{lstlisting}
// Fetch Cp(T) series from NIST thermochemical tables.
WebResponse fetch_cp_curve(const std::string& formula,
						   double T_min_K, double T_max_K);

// Fetch phase transition temperatures (Tm, Tb).
WebResponse fetch_phase_transitions(const std::string& formula);
\end{lstlisting}

Both go through \texttt{CachingFetcher} --- no raw network calls outside the
decorator.  Cache keys include formula, T\_min, T\_max rounded to 50 K bands.

\subsection{LocalSanityChecker extension (Tier 3)}

\begin{itemize}[leftmargin=*, itemsep=2pt]
  \item \texttt{thermo.melting\_point\_K > 0} if phase = solid
  \item \texttt{thermo.boiling\_point\_K > thermo.melting\_point\_K} if both present
  \item \texttt{cp\_T} curve has at least 2 points and is monotonically valid
		(no negative Cp)
  \item \texttt{delta\_Hf\_298} is in physically plausible range for the element
		family (actinides: $-1500$ to $+200$ kJ/mol; noble gas: always near 0)
  \item \texttt{eos.eos\_model} is non-empty
\end{itemize}

Failure tag: \texttt{thermo\_block\_invalid}.

\subsection{property\_values rows written by fill-thermo}

Thermo properties are stored as multiple rows, one per T/P point on the curve.
The \texttt{temperature\_K} column is the state variable; \texttt{pressure\_Pa}
defaults to 101325.0 for standard-state values.

\begin{center}
\begin{tabular}{lll}
\toprule
\textbf{property\_name} & \textbf{units} & \textbf{note} \\
\midrule
\texttt{delta\_Hf\_kJ\_mol}       & kJ/mol        & single row, T=298.15 \\
\texttt{delta\_Gf\_kJ\_mol}       & kJ/mol        & single row, T=298.15 \\
\texttt{S\_298\_J\_mol\_K}        & J/(mol K)     & single row, T=298.15 \\
\texttt{cp}                       & J/(mol K)     & one row per T point \\
\texttt{melting\_point\_K}        & K             & single row \\
\texttt{boiling\_point\_K}        & K             & single row \\
\texttt{critical\_T\_K}           & K             & single row \\
\texttt{critical\_P\_Pa}          & Pa            & single row \\
\texttt{acentric\_factor}         & dimensionless & single row \\
\texttt{density\_kg\_m3}          & kg/m\(^3\)    & one row per (T,P) point \\
\texttt{viscosity\_Pa\_s}         & Pa s          & one row per (T,P) point \\
\texttt{thermal\_conductivity\_W\_mK} & W/(m K)  & one row per (T,P) point \\
\bottomrule
\end{tabular}
\end{center}

\subsection{CLI sub-command}

\begin{lstlisting}[language={}]
vsepr-ufx ufx auto2 fill-thermo --batch 500 --db ufx_auto2.sqlite
vsepr-ufx ufx auto2 fill-thermo --batch 500 --nist-validate --db ufx_auto2.sqlite
\end{lstlisting}

\texttt{--nist-validate} fires NIST WebBook lookups for Tm, Tb, Cp(298) and writes
comparison rows to \texttt{validation\_records} with \texttt{validator\_id='nist\_thermo'}.

\subsection{Continual harness integration}

\begin{lstlisting}[language={}]
# Cycle order with Phase 7:
randomfill
fill-molecular
fill-thermo
validate       --batch $ValidateBatch
validate-web   --batch $WebBatch
promote        --min-score $MinScore
audit
\end{lstlisting}

\subsection{Audit success criteria after Phase 7}

\begin{itemize}[leftmargin=*, itemsep=2pt]
  \item \texttt{Total property values} grows by $\geq 12 \times$ \texttt{BatchSize} per cycle
		(12 scalar rows + curve points per record)
  \item \texttt{top rejected reasons} includes \texttt{thermo\_block\_invalid} for a
		non-zero fraction (confirms the check fires)
  \item NIST cross-check produces nonzero \texttt{web pass} and \texttt{web missing} counts
  \item No record reaches \texttt{validate-web} without thermo provenance rows
\end{itemize}

% =========================================================
\section{Phase 8 --- Crystal / Material Importer (Tier 4)}
% =========================================================

\begin{tierbox}
\textbf{Goal:} attach \texttt{CrystalBlock} data to solid-phase records using
Materials Project and OQMD as external references, lattice-parameter heuristics
as fallback, and space-group sampling as a new generation axis.
\end{tierbox}

\subsection{Scope}

Phase 8 applies only to records where \texttt{phase = solid}.  Gas and liquid
records skip \texttt{fill-crystal} without failure.  The generator checks
\texttt{rec.identity.phase == "solid"} before populating anything.

\subsection{Axis expansion}

\begin{itemize}[leftmargin=*, itemsep=2pt]
  \item \textbf{space\_group} --- sampled from per-element-family space group frequency tables
		(e.g., FCC metals favour Fm-3m; hexagonal metals favour P63/mmc)
  \item \textbf{lattice\_a\_angstrom} --- sampled from empirical radius-coordination correlations
  \item \textbf{packing\_fraction} --- sampled 0.52 (SC) to 0.74 (FCC/HCP)
  \item \textbf{energy\_above\_hull\_eV\_atom} --- sampled 0.0 to 0.25 (stability filter)
  \item \textbf{structure\_source} --- one of: \texttt{generated}, \texttt{mp\_lookup},
		\texttt{oqmd\_lookup}
\end{itemize}

\subsection{Generator class: CrystalGenerator}

File: \texttt{src/ufx\_auto2/crystal\_generator.cpp}

\begin{lstlisting}
class CrystalGenerator {
public:
	CrystalGenerator();

	// Populate rec.crystal from axis sample.
	// Phase: only runs for solid-phase records.
	// Returns false (no-op) if rec.identity.phase != "solid".
	bool fill(UFXMaterialRecord& rec, const AxisSample& s) const;

	// Estimate lattice parameter a from element radius + coordination.
	static double estimate_a_angstrom(int Z,
									  int coordination,
									  double ionic_radius_est_angstrom) noexcept;

	// Map coordination number + element family to most probable space group.
	static std::string probable_space_group(const std::string& family,
											int coordination) noexcept;

	// Compute unit cell volume from a, b, c, alpha, beta, gamma.
	static double unit_cell_volume(const Vec3& abc,
								   const Vec3& angles_deg) noexcept;
private:
	// Space group frequency table: family -> vector<(sg, weight)>
	static const std::vector<std::pair<std::string,double>>&
		sg_table_(const std::string& family) noexcept;
};
\end{lstlisting>

\subsection{Materials Project fetcher}

New file: \texttt{src/ufx\_auto2/materials\_project\_fetcher.cpp}

\begin{lstlisting}
class MaterialsProjectFetcher : public IWebFetcher {
public:
	// Requires MP_API_KEY environment variable.
	// Falls back to cache-only mode if key not set.
	explicit MaterialsProjectFetcher(const std::string& api_key = "");

	WebResponse fetch(const WebQuery& q) override;

private:
	std::string api_key_;
	std::string build_url_(const WebQuery& q) const;
	WebResponse parse_response_(const std::string& body,
								const WebQuery& q) const;
};
\end{lstlisting}

\begin{rulebox}
If \texttt{MP\_API\_KEY} is not set, \texttt{MaterialsProjectFetcher} operates in
cache-read-only mode.  It will use existing cached responses but not make new
HTTP requests.  This is the correct offline behavior; it is not an error.
\end{rulebox}

The fetcher queries:
\texttt{https://api.materialsproject.org/materials/summary/?formula=\{formula\}\&fields=structure,formation\_energy\_per\_atom,energy\_above\_hull,band\_gap}

Results are cached under \texttt{web\_cache/mp/\{sha256\_key\}.json}.

\subsection{LocalSanityChecker extension (Tier 4)}

Applied only when \texttt{rec.identity.phase == "solid"}:

\begin{itemize}[leftmargin=*, itemsep=2pt]
  \item \texttt{crystal.space\_group} is non-empty
  \item \texttt{crystal.lattice\_abc\_angstrom.x > 0}
  \item \texttt{crystal.density\_g\_cm3 > 0}
  \item \texttt{crystal.packing\_fraction} between 0.10 and 0.82
  \item \texttt{crystal.energy\_above\_hull\_eV\_atom} between $-0.1$ and $1.5$
		(large negative values indicate an error in the estimator)
\end{itemize}

Failure tag: \texttt{crystal\_block\_invalid}.

\subsection{property\_values rows written by fill-crystal}

\begin{center}
\begin{tabular}{lll}
\toprule
\textbf{property\_name} & \textbf{units} & \textbf{note} \\
\midrule
\texttt{space\_group}              & string    & stored as \texttt{value\_text} \\
\texttt{space\_group\_number}      & ITA int   & stored as \texttt{value\_real} \\
\texttt{lattice\_a}                & \AA{}     & \\
\texttt{lattice\_b}                & \AA{}     & \\
\texttt{lattice\_c}                & \AA{}     & \\
\texttt{lattice\_alpha\_deg}       & degrees   & \\
\texttt{unit\_cell\_volume}        & \AA{}$^3$ & \\
\texttt{density\_g\_cm3}           & g/cm$^3$  & at T=298.15, P=101325 \\
\texttt{packing\_fraction}         & \textemdash & \\
\texttt{formation\_energy\_eV\_atom} & eV/atom  & \\
\texttt{energy\_above\_hull\_eV\_atom} & eV/atom & \\
\texttt{band\_gap\_eV}             & eV        & \texttt{NaN} if metallic \\
\texttt{debye\_temperature\_K}     & K         & estimated from elastic modulus \\
\bottomrule
\end{tabular}
\end{center}

\subsection{CLI sub-command}

\begin{lstlisting}[language={}]
vsepr-ufx ufx auto2 fill-crystal --batch 500 --db ufx_auto2.sqlite
vsepr-ufx ufx auto2 fill-crystal --batch 500 --mp-validate --db ufx_auto2.sqlite
\end{lstlisting}

\texttt{--mp-validate} fires Materials Project API lookups for formation energy and
space group.  Writes comparison rows to \texttt{validation\_records} with
\texttt{validator\_id='materials\_project'}.

\subsection{Continual harness integration}

\begin{lstlisting}[language={}]
# Cycle order with Phase 8:
randomfill
fill-molecular
fill-thermo
fill-crystal
validate       --batch $ValidateBatch
validate-web   --batch $WebBatch
promote        --min-score $MinScore
audit
\end{lstlisting>

\subsection{Audit success criteria after Phase 8}

\begin{itemize}[leftmargin=*, itemsep=2pt]
  \item Solid-phase records have nonzero \texttt{property\_values} rows with
		\texttt{block\_name='crystal'}
  \item Gas/liquid records have zero crystal rows (phase filter working)
  \item If MP API key is set: nonzero \texttt{validator\_id='materials\_project'} rows
  \item \texttt{crystal\_block\_invalid} appears in top rejected reasons
\end{itemize}

% =========================================================
\section{Phase 9 --- Mechanics + Transport Macro Bridge (Tier 5)}
% =========================================================

\begin{tierbox}
\textbf{Goal:} populate \texttt{MechanicalBlock} and \texttt{TransportBlock} using
atomistic-to-engineering bridge correlations, and connect them to the macro
observable layer needed for pipeline erosion, thermal coupling, and flow modeling.
\end{tierbox}

\subsection{Context: the atomistic-to-macro chain}

Phases 2--8 built atomic-scale and molecular properties.  Phase 9 generates
the engineering-scale observables that are actually used in macro simulations:

\begin{center}
\texttt{UFF curvature $\rightarrow$ $E_\text{eff}$}
\quad$\cdot$\quad
\texttt{Thermo $\rightarrow$ $k_\text{eff}$}
\quad$\cdot$\quad
\texttt{Crystal density $\rightarrow$ bulk modulus estimate}
\quad$\cdot$\quad
\texttt{EOS viscosity $\rightarrow$ Darcy coefficient}
\end{center}

These are not raw DFT results.  They are engineering estimates with explicit
approximation labels.  Label them.

\subsection{Axis expansion}

\begin{itemize}[leftmargin=*, itemsep=2pt]
  \item \textbf{strain\_rate} --- sampled $10^{-6}$ to $10^3$ s$^{-1}$
  \item \textbf{test\_environment} --- one of: \texttt{vacuum}, \texttt{air},
		\texttt{molten\_salt}, \texttt{brine}, \texttt{hydrogen}
  \item \textbf{porosity} --- sampled 0.0 to 0.60
  \item \textbf{tortuosity} --- sampled 1.0 to 5.0
  \item \textbf{thermal\_gradient\_K\_m} --- sampled 0 to $10^4$
  \item \textbf{flow\_velocity\_m\_s} --- sampled 0.0 to 100.0
\end{itemize}

\subsection{Generator class: MacroPropertyGenerator}

File: \texttt{src/ufx\_auto2/macro\_property\_generator.cpp}

\begin{lstlisting}
class MacroPropertyGenerator {
public:
	MacroPropertyGenerator();

	// Fill rec.mechanical using UFF stiffness-to-modulus bridge.
	// Approximation: E ~ (r1 / D1) * coordination * k_harmonic_factor
	void fill_mechanical(UFXMaterialRecord& rec,
						 const AxisSample& s) const;

	// Fill rec.transport using Bridgman correlations from thermo data.
	// k_eff from thermal conductivity model (kinetic theory + phonon scattering)
	// diffusivity from Stokes-Einstein / empirical scaling
	void fill_transport(UFXMaterialRecord& rec,
						const AxisSample& s) const;

	// Effective thermal conductivity via Maxwell-Garnett (porous medium).
	static double k_eff_maxwell_garnett(double k_solid,
										double k_fluid,
										double porosity) noexcept;

	// Effective Young's modulus via Voigt bound.
	static double E_eff_voigt(double E_bulk, double porosity) noexcept;

	// Darcy permeability from Kozeny-Carman equation.
	static double permeability_kozeny_carman(double d_particle_m,
											 double porosity) noexcept;

private:
	static double k_phonon_estimate_(int Z, int period,
									 double debye_T_K) noexcept;
};
\end{lstlisting>

\subsection{Bridge correlations used in Phase 9}

\begin{center}
\begin{tabular}{lll}
\toprule
\textbf{Output property} & \textbf{Input} & \textbf{Model} \\
\midrule
$E_\text{eff}$ (GPa)          & $r_1$, $D_1$, coord, density   & UFF stiffness bridge \\
$G_\text{eff}$ (GPa)          & $E_\text{eff}$, Poisson est.    & isotropic elasticity \\
$K_\text{eff}$ (GPa)          & $E_\text{eff}$, Poisson est.    & isotropic elasticity \\
$k_\text{eff}$ (W/m\,K)       & Debye T, density, $C_p$        & Bridgman / kinetic theory \\
$D$ (m$^2$/s)                 & mol.\ weight, $T$, viscosity    & Stokes-Einstein \\
Darcy permeability (m$^2$)    & $d_\text{particle}$, porosity   & Kozeny-Carman \\
\bottomrule
\end{tabular}
\end{center}

\textbf{Every output carries:}
\texttt{method\_tag = "phase9\_bridge\_heuristic"},
\texttt{confidence = 0.30}, \texttt{approximation\_label} set.

\subsection{LocalSanityChecker extension (Tier 5)}

\begin{itemize}[leftmargin=*, itemsep=2pt]
  \item \texttt{mechanical.E\_eff\_GPa} is in range $[0.001, 1200]$ GPa
  \item \texttt{transport.k\_eff\_W\_mK} is positive
  \item \texttt{transport.diffusivity\_m2\_s} is positive and below $10^{-4}$
  \item \texttt{transport.porosity} between 0.0 and 0.98 if present
  \item \texttt{mechanical.approximation\_label} non-empty
		(enforces the ``label the approximation'' hard rule)
\end{itemize}

Failure tag: \texttt{macro\_block\_invalid}.

\subsection{property\_values rows written by fill-transport and fill-mechanical}

\begin{center}
\begin{tabular}{lll}
\toprule
\textbf{property\_name} & \textbf{units} & \textbf{block\_name} \\
\midrule
\texttt{E\_eff\_GPa}              & GPa        & mechanical \\
\texttt{G\_eff\_GPa}              & GPa        & mechanical \\
\texttt{K\_eff\_GPa}              & GPa        & mechanical \\
\texttt{poisson\_ratio}           & \textemdash & mechanical \\
\texttt{yield\_MPa}               & MPa        & mechanical \\
\texttt{k\_eff\_W\_mK}            & W/(m K)    & transport \\
\texttt{diffusivity\_m2\_s}       & m$^2$/s    & transport \\
\texttt{permeability\_m2}         & m$^2$      & transport \\
\texttt{porosity}                 & \textemdash & transport \\
\texttt{tortuosity}               & \textemdash & transport \\
\texttt{darcy\_coeff}             & \textemdash & transport \\
\texttt{mass\_transfer\_coeff}    & m/s        & transport \\
\bottomrule
\end{tabular}
\end{center}

\subsection{CLI sub-commands}

\begin{lstlisting}[language={}]
vsepr-ufx ufx auto2 fill-transport  --batch 500 --db ufx_auto2.sqlite
vsepr-ufx ufx auto2 fill-mechanical --batch 500 --db ufx_auto2.sqlite
\end{lstlisting}

These may be called as separate commands (for partial fills) or as a combined
\texttt{fill-macro} alias that calls both.

\subsection{Continual harness integration}

\begin{lstlisting}[language={}]
# Cycle order with Phase 9 (full):
randomfill
fill-molecular
fill-thermo
fill-crystal
fill-transport
fill-mechanical
validate       --batch $ValidateBatch
validate-web   --batch $WebBatch
promote        --min-score $MinScore
audit
\end{lstlisting}

\begin{phasebox}
\textbf{Performance note:} each fill step should process only records that are
missing the target block rows.  The query pattern is identical to the
\texttt{validate} back-fill query:
\begin{lstlisting}[language=SQL]
SELECT id FROM material_records
WHERE NOT EXISTS (
  SELECT 1 FROM property_values pv
  WHERE pv.material_id = material_records.id
  AND   pv.block_name  = 'transport'
)
LIMIT ?;
\end{lstlisting}
This is idempotent.  Cycles do not re-fill already-filled records.
\end{phasebox}

\subsection{Audit success criteria after Phase 9}

\begin{itemize}[leftmargin=*, itemsep=2pt]
  \item \texttt{Total property values} $\geq 30 \times$ record count
  \item \texttt{mechanical} and \texttt{transport} block rows present for all
		non-rejected records
  \item \texttt{macro\_block\_invalid} in top rejected reasons for out-of-range
		generated estimates
  \item Average composite score rises above 0.85 (macro bridge contributes
		to $S_\text{macro}$ component)
\end{itemize}

% =========================================================
\section{Phase 10 --- Meta-Score Active Scheduler (Tier 9)}
% =========================================================

\begin{tierbox}
\textbf{Goal:} compute \texttt{MetaScoreBlock} fields for every promoted record,
then use them to bias the axis sampler toward under-covered, high-usefulness,
low-risk regions of parameter space.  Stop generating the same octahedral iron
oxide at 298 K for the thousandth time.
\end{tierbox}

\subsection{What Phase 10 is and is not}

Phase 10 is \textbf{not} a machine learning system.  It is a deterministic
coverage-aware scheduler built on the \texttt{CoverageMap} that already exists
in Phase 2.

Phase 10 extends \texttt{CoverageMap} with:
\begin{itemize}[leftmargin=*, itemsep=2pt]
  \item per-region \texttt{coverage\_score} (record count / region capacity)
  \item per-region \texttt{weirdness\_avg} (mean weirdness of existing records in region)
  \item per-region \texttt{macro\_usefulness\_avg} (mean macro usefulness)
  \item per-region \texttt{extrapolation\_risk} (fraction of records in region with
		confidence $< 0.5$)
\end{itemize}

\subsection{Generator class: MetaScoreEngine}

File: \texttt{src/ufx\_auto2/meta\_score\_engine.cpp}

\begin{lstlisting}
class MetaScoreEngine {
public:
	MetaScoreEngine();

	// Compute all MetaScoreBlock fields for a fully-filled record.
	// Requires: identity, force_field, molecular, thermo, and crystal
	// blocks populated (Phases 2-8).
	void compute(UFXMaterialRecord& rec,
				 const CoverageMap& coverage) const;

	// Bias weight for choosing the next weak axis region.
	// Higher = more likely to be sampled next.
	double region_priority(const RegionKey& region,
						   const CoverageMap& coverage) const noexcept;

	// Update coverage map meta fields after a record is promoted.
	void update_coverage(CoverageMap& coverage,
						 const UFXMaterialRecord& rec) const;

private:
	double weirdness_(const UFXMaterialRecord& rec,
					  const CoverageMap& coverage) const noexcept;
	double novelty_   (const UFXMaterialRecord& rec,
					  const CoverageMap& coverage) const noexcept;
	double macro_usefulness_(const UFXMaterialRecord& rec) const noexcept;
	double extrapolation_risk_(const UFXMaterialRecord& rec) const noexcept;
};
\end{lstlisting>

\subsection{Composite score model}

From \texttt{UFX\_continual\_2.tex} \S19:

\begin{equation*}
S = 0.25 S_\text{identity} + 0.20 S_\text{local} + 0.20 S_\text{reference} + 0.15 S_\text{web} + 0.10 S_\text{repeat} + 0.10 S_\text{macro}
\end{equation*}

Phase 10 populates the $S_\text{macro}$ component for the first time.
Prior phases contribute to the others.  The composite score is stored in
\texttt{property\_values} with \texttt{block\_name='meta'} and
\texttt{property\_name='composite\_score'}.

\subsection{Sampling bias: weak-region priority}

The \texttt{RandomAxisSampler} accepts an optional priority weight map.
Phase 10 injects this map at the start of each cycle via
\texttt{CoverageMap::region\_weights()}:

\begin{lstlisting}
// Priority function:
// p(region) = (1 - coverage_score)^2
//            * (1 + macro_usefulness_avg)
//            * (1 - extrapolation_risk)
//            * (1 + weirdness_avg * 0.2)
//
// Normalization: weights are renormalized to sum to 1.0 before sampling.
std::map<RegionKey, double> CoverageMap::region_weights() const;
\end{lstlisting}

This keeps the sampler deterministic (given the same seed and coverage state)
while steering generation away from saturated, low-usefulness regions.

\subsection{MetaScoreBlock writes to property\_values}

\begin{center}
\begin{tabular}{lll}
\toprule
\textbf{property\_name} & \textbf{range} & \textbf{note} \\
\midrule
\texttt{weirdness}              & 0.0--1.0 & vs.\ periodic neighbours \\
\texttt{coverage}               & 0.0--1.0 & region record density \\
\texttt{confidence}             & 0.0--1.0 & aggregate validation \\
\texttt{novelty}                & 0.0--1.0 & distance from validated entries \\
\texttt{local\_consistency}     & 0.0--1.0 & agreement with neighbours \\
\texttt{source\_conflict}       & 0.0--1.0 & inter-source disagreement \\
\texttt{macro\_usefulness}      & 0.0--1.0 & predicts $E$, $k$, dP, damage \\
\texttt{simulation\_stability}  & 0.0--1.0 & VSEPR run stability \\
\texttt{extrapolation\_risk}    & 0.0--1.0 & model outside valid space \\
\texttt{composite\_score}       & 0.0--1.0 & promotion gate value \\
\bottomrule
\end{tabular}
\end{center}

\subsection{CLI sub-command}

\begin{lstlisting}[language={}]
vsepr-ufx ufx auto2 score --batch 500 --db ufx_auto2.sqlite
vsepr-ufx ufx auto2 score --batch 500 --recompute --db ufx_auto2.sqlite
\end{lstlisting}

\texttt{--recompute} recalculates meta scores for all promoted records, not just
those missing scores.  Use after coverage map changes significantly.

\subsection{Audit extensions for Phase 10}

The \texttt{ufx auto2 audit} command is extended to show:

\begin{lstlisting}[language={}]
  -- Meta-Score Summary --
  Avg composite score   : 0.8723
  Highest weirdness     : Lu_ox3_c8_cubic_solid (0.91)
  Highest macro use     : Fe_ox2_c6_octahedral_solid (0.87)
  Highest extrap risk   : Og_ox0_c1_linear_gas (0.95)
  Weakest region        : actinide::molten_salt::octahedral
  Priority bias active  : yes  (Phase 10 scheduler)
\end{lstlisting}

\subsection{Continual harness integration}

\begin{lstlisting}[language={}]
# Cycle order with Phase 10 (full pipeline):
randomfill     --count $BatchSize
fill-molecular --batch $BatchSize
fill-thermo    --batch $BatchSize
fill-crystal   --batch $BatchSize
fill-transport --batch $BatchSize
fill-mechanical --batch $BatchSize
score          --batch $BatchSize
validate       --batch $ValidateBatch
validate-web   --batch $WebBatch
promote        --min-score $MinScore
audit
\end{lstlisting>

\begin{phasebox}
The harness parameter \texttt{-SleepSeconds} absorbs any latency from the
extended fill chain.  For the 1 GB overnight run, \texttt{-SleepSeconds 3} is
the minimum.  For the 10 GB run, use \texttt{-SleepSeconds 5} to avoid
overwhelming the web cache rate limiter.
\end{phasebox}

% =========================================================
\section{Schema Additions}
% =========================================================

Phases 6--10 add no new tables.  They use the existing generic
\texttt{property\_values} table with correct \texttt{block\_tier} and
\texttt{block\_name} values.  Dedicated tables for thermo curves, crystal
structures, or mechanical tensors may be added post-Phase 10 if query patterns
demand it, per the hard rule in \S\ref{sec:hardrules}.

The one schema extension is a \textbf{meta-score summary view} for audit
performance:

\begin{lstlisting}[language=SQL]
CREATE VIEW IF NOT EXISTS v_meta_scores AS
SELECT
	mr.id,
	mr.material_key,
	mr.source_class,
	MAX(CASE WHEN pv.property_name = 'composite_score'
			 THEN pv.value_real END)  AS composite_score,
	MAX(CASE WHEN pv.property_name = 'weirdness'
			 THEN pv.value_real END)  AS weirdness,
	MAX(CASE WHEN pv.property_name = 'macro_usefulness'
			 THEN pv.value_real END)  AS macro_usefulness,
	MAX(CASE WHEN pv.property_name = 'extrapolation_risk'
			 THEN pv.value_real END)  AS extrapolation_risk
FROM material_records mr
JOIN property_values pv ON pv.material_id = mr.id
WHERE pv.block_name = 'meta'
GROUP BY mr.id;
\end{lstlisting}

This view is created by \texttt{ufx\_auto2\_init\_db()} alongside the existing
six tables.

% =========================================================
\section{CMake Integration}
% =========================================================

Each phase adds one or two source files to the \texttt{vsepr-ufx} target in
\texttt{cmake/CoreBuild.cmake}.  The pattern is identical to the Phase 5 web
source additions.

\begin{lstlisting}[language={}]
# Phase 6
src/ufx_auto2/molecular_descriptor_generator.cpp

# Phase 7
src/ufx_auto2/thermo_generator.cpp
src/ufx_auto2/nist_thermo_fetcher.cpp    # extends NISTFetcher with Cp query

# Phase 8
src/ufx_auto2/crystal_generator.cpp
src/ufx_auto2/materials_project_fetcher.cpp

# Phase 9
src/ufx_auto2/macro_property_generator.cpp

# Phase 10
src/ufx_auto2/meta_score_engine.cpp
\end{lstlisting}

No new link dependencies are needed beyond those already in the Phase 5 build
(\texttt{sqlite3}, \texttt{winhttp}).  The Materials Project fetcher uses the
same WinHTTP path as PubChem and NIST.

% =========================================================
\section{File Layout Summary}
% =========================================================

\begin{center}
\begin{tabular}{ll}
\toprule
\textbf{File} & \textbf{Phase} \\
\midrule
\texttt{include/ufx\_auto2/smiles\_vocab.hpp}                   & 6 \\
\texttt{src/ufx\_auto2/molecular\_descriptor\_generator.cpp}     & 6 \\
\texttt{src/ufx\_auto2/thermo\_generator.cpp}                    & 7 \\
\texttt{src/ufx\_auto2/nist\_thermo\_fetcher.cpp}                & 7 \\
\texttt{src/ufx\_auto2/crystal\_generator.cpp}                   & 8 \\
\texttt{src/ufx\_auto2/materials\_project\_fetcher.cpp}          & 8 \\
\texttt{include/ufx\_auto2/materials\_project\_fetcher.hpp}      & 8 \\
\texttt{src/ufx\_auto2/macro\_property\_generator.cpp}           & 9 \\
\texttt{src/ufx\_auto2/meta\_score\_engine.cpp}                  & 10 \\
\texttt{include/ufx\_auto2/meta\_score\_engine.hpp}              & 10 \\
\midrule
\texttt{src/cli/cmd\_ufx.cpp}  & modified (6--10 sub-commands) \\
\texttt{src/cli/cmd\_ufx.hpp}  & modified (6--10 handler declarations) \\
\texttt{src/v4/uff/ufx\_schema.cpp} & modified (Phase 10 view, score queries) \\
\texttt{cmake/CoreBuild.cmake} & modified (new sources added) \\
\texttt{tools/run\_ufx\_auto2\_continual.ps1} & modified (fill steps in cycle) \\
\bottomrule
\end{tabular}
\end{center}

% =========================================================
\section{Incremental Run Order}
% =========================================================

Do not implement all five phases simultaneously.  The correct order:

\begin{enumerate}[leftmargin=*, itemsep=4pt]
  \item \textbf{Implement Phase 6 only.}  Run medium test (250 MB target).
		Confirm \texttt{property\_values} rows appear, molecular weight
		cross-check fires, audit shows nonzero Tier 2 data.

  \item \textbf{Implement Phase 7 only.}  Run medium test again.
		Confirm Cp curves appear in \texttt{property\_values}, NIST thermo
		validator produces nonzero pass/missing counts.

  \item \textbf{Implement Phase 8 only.}  Run 1 GB overnight candidate.
		Confirm solid-phase records gain crystal block, gas/liquid records
		do not.  MP API either fires (if key set) or operates in cache-read mode.

  \item \textbf{Implement Phase 9 only.}  Run 1 GB overnight candidate.
		Confirm transport and mechanical rows appear.  Average composite score
		should rise visibly.

  \item \textbf{Implement Phase 10.}  Run the 10 GB candidate.
		Confirm the weakest-region audit entry changes each cycle
		(scheduler is steering).  Weirdness and macro-usefulness scores appear
		in audit output.
\end{enumerate}

\begin{rulebox}
Run the tiny test (25 MB, \texttt{-EmergencyStopOnFail}) after each phase
before committing to the longer runs.  This is the lowest-cost proof that
command wiring is correct and the harness cycle does not deadlock.
\end{rulebox}

% =========================================================
\section{Hard Rules}
\label{sec:hardrules}
% =========================================================

\begin{rulebox}
\begin{enumerate}[leftmargin=*, itemsep=6pt]

  \item \textbf{Fill steps process only records missing the target block.}
		The missing-block query is idempotent.  Rerunning \texttt{fill-thermo}
		on a database where thermo rows already exist is a no-op.  This is
		correct behavior and not a performance regression.

  \item \textbf{Every generated value carries an approximation label.}
		Shomate heuristics, Bridgman correlations, UFF bridge estimates ---
		all must have \texttt{method\_tag} and \texttt{confidence $\leq$ 0.40}
		unless validated by an external source.

  \item \textbf{State columns are not optional.}  \texttt{temperature\_K},
		\texttt{pressure\_Pa}, and \texttt{phase} must be populated on every
		\texttt{property\_values} row.  The generation fallback is T=298.15,
		P=101325.0, phase=identity.phase.

  \item \textbf{Generic \texttt{property\_values} before dedicated tables.}
		Do not create \texttt{thermo\_properties} or \texttt{crystal\_properties}
		tables until the query patterns from the 1 GB run clearly require them.

  \item \textbf{Phase 10 is a scheduler, not a filter.}  It must not silently
		drop records from the cycle.  It changes what gets generated next,
		not what gets validated or promoted.

  \item \textbf{Nuclear/radiation fields never enter these fill steps.}
		Tier 7 (\texttt{RadiationBlock}) follows after Phase 10 is stable.
		The isotope/decay axis exists in \texttt{UFXMaterialRecord} and
		\texttt{AxisSample} but is not sampled until Tier 1 through 5 are stable.

  \item \textbf{The harness cycle order is fixed.}  Fill steps run before
		validate.  Score (\texttt{fill-meta}) runs before promote.
		Promote never runs before validate-web on a \texttt{-WebValidate} run.

  \item \textbf{SQLite is truth.  Audit is the instrument panel.}  If the
		audit output does not show the expected block counts after a fill step,
		the fill step is broken --- not the audit.

\end{enumerate}
\end{rulebox}

\vspace{10pt}
\hrule
\vspace{8pt}

\begin{center}
\textit{Phase 2 proved the skeleton survives.\\
Phases 6--10 teach it chemistry, thermodynamics, structure, mechanics, and taste.\\
In that order.  Not all at once.}
\end{center}

\end{document}
